\chapter{اللقاء السابع: بداية سورة البقرة وتأويل فواتح السور}

\section{جواز تسمية سورة البقرة}
السلام عليكم ورحمة الله وبركاته. الحمد لله رب العالمين، والصلاة والسلام على نبينا محمد وعلى آله وصحبه أجمعين. نبدأ بإذن الله تبارك وتعالى في اللقاء السابع من دراستنا لتفسير الإمام \rawi{ابن جرير الطبري} رحمه الله، وقد وصلنا إلى بداية سورة البقرة.

\isnad{قال الطبري رحمه الله:} «القول في تفسير السورة التي يذكر فيها البقرة».
لعل \rawi{الطبري} اعتمد في هذه التسمية وتجنب قول "سورة البقرة" على حديث مروي \isnad{عن} \rawi{أنس} رضي الله عنه عن النبي صلى الله عليه وسلم قال: \begin{hadith}لا تقولوا سورة البقرة ولا سورة آل عمران... ولكن قولوا السورة التي يذكر فيها البقرة\end{hadith}.

وأود التنبيه على أن هذا الحديث \textbf{منكر} لا يثبت؛ فقد حكم عليه بالنكارة الإمام \rawi{أحمد بن حنبل}، وكذلك أنكره \rawi{البيهقي} في كتابه، وضعفه \rawi{ابن كثير} رحمه الله. 
وقد بوّب الإمام \rawi{البخاري} في صحيحه: «باب من لم ير بأساً أن يقول سورة البقرة وسورة كذا»، وذكر أحاديث صحيحة ترد على هذا، منها حديث \rawi{ابن مسعود} رضي الله عنه حين رمى جمرة العقبة فقال: \begin{hadith}من هاهنا والذي لا إله غيره رمى الذي أُنزلت عليه سورة البقرة صلى الله عليه وسلم\end{hadith}. فهذا الحديث الصحيح فيه رد على من يكره تسميتها بسورة البقرة.

\separator

\section{تأويل (الم) وفواتح السور}
\isnad{قال أبو جعفر:} «القول في تأويل قوله جل ثناؤه: \begin{ayah}الم\end{ayah}. اختلف تراجمة القرآن في تأويل قول الله تعالى ذكره \textit{الم}. فقال بعضهم: هي اسم من أسماء القرآن... وقال بعضهم: هي فواتح يفتتح الله بها القرآن... وقال بعضهم: هي اسم للسورة... وقال بعضهم: هو اسم الله الأعظم... وقال بعضهم: هو قسم أقسم الله به... وقال بعضهم: هي حروف مقطعة من أسماء وأفعال، كل حرف من ذلك بمعنى غير معنى الحرف الآخر...».

\begin{fawaid}[فواتح السور مما يعلم معناه]
مجرد كلام الصحابة والتابعين في الأحرف التي في فواتح السور، واجتهادهم في بيان معناها وإن اختلفوا؛ يدل دلالة واضحة على أنها ليست من القسم المتشابه الذي لا يعلمه إلا الله (كما يزعم البعض). فمن قال: "لا علم لي بها" فقد حكم على نفسه، ولكن الادعاء بأنها مما يختص الله بعلمه فهذا محل نقد.
\end{fawaid}

\subsection{منهج الطبري في عرض الأقوال والرد عليها}
ساق \rawi{الطبري} الأقوال الواردة عن \rawi{ابن عباس} و\rawi{مجاهد} و\rawi{السدي} وغيرهم، ثم ساق أقوال أهل العربية. ومن الملاحظ في منهجه في هذا الموضع:

\begin{enumerate}
    \item \textbf{استبعاد الرواية شديدة الضعف:} ذكر \rawi{الطبري} قولاً بأن هذه الحروف من "حساب الجُمَّل" (وهو ضرب من الحساب يجعل لكل حرف عدداً)، ولكنه قال: «كرهنا ذكر الذي حُكي ذلك عنه إذ كان الذي رواه ممن لا يُعتمد على روايته ونقله». فإذا تحقق لديه ضعف الرواية وأنها لا تُعتمد، فإنه يُعرض عن ذكرها وسندها.
    \item \textbf{تأخير الرد والتفنيد:} قد يذكر \rawi{الطبري} القول ويذكر حجته ووجهه، ثم يؤخر الرد عليه صفحات طويلة. فلا ينبغي لقارئ الطبري أن يتعجل ويظن أنه ما دام قد ساق القول ووجهه فإنه يرتضيه، بل عادته أن يجمع الأقوال وحججها، ثم يشرع في تفنيدها في آخر المسألة.
    \item \textbf{أسلوب الإلزام في المناظرة:} استعمل \rawi{الطبري} "الإلزام" لرد الأقوال المخالفة. فمثلاً، اعترض معترض: كيف يكون الحرف الواحد دالاً على معانٍ كثيرة مختلفة؟ فألزمه الطبري بأن الكلمة الواحدة في لسان العرب (مثل: أمة، ودين) تدل على معانٍ مختلفة، فإن أقرّ بذلك في الكلمات لزمه الإقرار به في الحروف.
\end{enumerate}

وهنا نبه الشيخ إلى فرق دقيق: كون هذه السور تفتتح بحروف مقطعة أمر لا نزاع فيه من جهة الوصف الشكلي، لكن جعله نفس الجواب التفسيري ليس هو مراد كل من تكلم فيها. فالوصف بأنها "حروف مقطعة" شيء، والقول في معناها ودلالتها شيء آخر.

\subsection{خلاصة رأي الطبري في الأحرف المقطعة}
\isnad{قال أبو جعفر:} «والصواب عندي من القول في تأويل مفاتح السور التي هي حروف المعجم، أن الله جل ثناؤه جعلها حروفاً مقطعة ولم يصل بعضها ببعض... لأنه عز ذكره أراد بلطفه الدلالة بكل حرف منه على معان كثيرة لا على معنى واحد... فكل حرف منه يحوي ما قاله \rawi{الربيع} وما قاله سائر المفسرين غيره فيه».

فالطبري يرجح أن الحرف الواحد يشتمل على دلالات متعددة (أسماء لله، آلاء، أزمنة)، وخطّأ قول اللغويين الذين زعموا أنها مجرد حروف هجاء استُغني بذكر بعضها عن باقيها، محتجاً بأن قولهم هذا خروج عن إجماع الصحابة والتابعين، و«كفى دلالة على خطئه شهادة الحجة عليه بالخطأ».

\separator

\section{تأويل (ذلك الكتاب لا ريب فيه)}
\isnad{قال أبو جعفر:} «القول في تأويل قوله جل ثناؤه: \begin{ayah}ذَٰلِكَ الْكِتَابُ\end{ayah}. قال عامة المفسرين: تأويل قول الله: \textit{ذلك الكتاب}، هذا الكتاب».

ثم طرح إشكالاً: كيف يجوز أن يكون (ذلك) - وهي إشارة للغائب - بمعنى (هذا) - وهي إشارة للحاضر؟
وأجاب بأن ما تقضّى وقرب تقضيه من الأخبار يُنزل منزلة الحاضر، كقولك لمن حدثك بحديث: "ذلك والله كما قلت"، أو "هذا والله كما قلت". فلما تقدم ذكر (الم)، صار الخبر عنه لقرب تقضيه كالحاضر المشاهد.

\begin{fawaid}[منهج الطبري في بيان لوازم الأقوال]
من دقة الطبري أنه يذكر ما يترتب على القول؛ فبعد أن وجّه مجيء (ذلك) بمعنى (هذا)، ذكر قولاً آخر بأن المراد بالكتاب هنا: التوراة والإنجيل، وعلّق قائلاً: «وإذا وُجِّه تأويل ذلك إلى هذا الوجه، فلا مؤونة (أي لا تكلف) فيه على متأوله»؛ لأن الإشارة حينئذ ستكون لغائب فعلاً على حقيقتها دون تأويل.
\end{fawaid}

\begin{fawaid}[نقل اتفاق المفسرين له وزنه في الترجيح]
قول \rawi{الطبري} هنا: «قال عامة المفسرين» ليس عبارة عابرة، بل يدل على اعتبار وزن اتفاق أهل التفسير في الباب، وأنه قد يجعل هذا الأصل منطلقاً ثم يذكر ما يَرِد عليه من الإشكال ويوجهه، لا أن يبدأ دائماً من فراغ.
\end{fawaid}

\isnad{وأما قوله:} \begin{ayah}لَا رَيْبَ ۛ فِيهِ\end{ayah} أي: لا شك فيه.

\separator

\section{تأويل (هدى للمتقين)}
\isnad{قال أبو جعفر:} «هدى من الضلالة... نور للمتقين... ولو كان نوراً لغير المتقين ورشاداً لغير المؤمنين، لم يخصص الله عز وجل المتقين بأنه لهم هدى، بل كان يعم به جميع المنذرين... فالمؤمن به مهتد والكافر به محجوب».

\begin{fawaid}[الفرق بين هداية البيان وهداية الانتفاع]
قد يُشكل كلام الطبري هنا لمن ظن أنه يحصر كون القرآن هدى للمتقين فقط. والصواب أن يقال: القرآن (هدى للناس) من جهة الإرشاد والبيان والإنذار، ولكنه (هدى للمتقين) من جهة الانتفاع والتوفيق. فالغيث ينزل على كل الأرض، ولكن لا تنتفع به وتنبت إلا الأرض الطيبة.
\end{fawaid}

\separator

\section{تأويل (الذين يؤمنون بالغيب)}
\isnad{قال أبو جعفر:} «ومعنى الإيمان عند العرب: التصديق، فيدعى المصدق بالشيء قولا مؤمناً به، ويدعى المصدق قوله بفعله مؤمناً... وقد تدخل الخشية لله في معنى الإيمان... والإيمان كلمة جامعة للإقرار بالله وكتبه ورسله، وتصديق الإقرار بالفعل... فالذي هو أولى بتأويل الآية... أن يكونوا موصوفين بالتصديق بالغيب قولا واعتقادا وعملا».

\begin{fawaid}[تعقيب على تفسير الإيمان بالتصديق]
تفسير الطبري للإيمان لغوياً بـ (التصديق) هو موضع منازعة؛ وقد اعتنى شيخ الإسلام \rawi{ابن تيمية} ببيان أن لفظ "الإيمان" يفترق عن "التصديق". فالإيمان فيه معنى (الائتمان) والمتابعة والانقياد، ولا يكون إلا في الإخبار بالغيب، بخلاف التصديق الذي يكون للغيب والشهادة. ولذا يُحمد للطبري أنه رغم تفسيره اللغوي، أكد شرعاً أن الإيمان كلمة جامعة تشمل القول والعمل، راداً بذلك على غلاة المرجئة والجهمية.
\end{fawaid}

\isnad{وأما (الغيب):} فهو ما غاب عن العباد من أمر الجنة والنار، وما أخبر به النبي صلى الله عليه وسلم ولم يكن لهم به علم من قبل.

\separator

\section{من المراد بهذه الآيات الأولى من سورة البقرة؟}
اختلف المفسرون فيمن نزلت فيهم هذه الآيات الأربع. 
ذكر \rawi{الطبري} الأقوال:
\begin{itemize}
    \item أنها في مؤمني العرب خاصة.
    \item أنها في مؤمني أهل الكتاب خاصة.
    \item أنها تعم جميع المؤمنين من العرب والعجم وأهل الكتاب.
\end{itemize}
ورجح \rawi{الطبري} القول بأنها مقسمة: فالآيتان الأوليان \begin{ayah}الَّذِينَ يُؤْمِنُونَ بِالْغَيْبِ\end{ayah} نزلتا في مؤمني العرب الذين لم يكن لهم كتاب من قبل. والآيتان اللتان تليهما \begin{ayah}وَالَّذِينَ يُؤْمِنُونَ بِمَا أُنزِلَ إِلَيْكَ وَمَا أُنزِلَ مِن قَبْلِكَ\end{ayah} نزلتا في مؤمني أهل الكتاب. واستدل بأن الله صنف الكفار صنفين (كافر صريح ومنافق)، فكذلك صنف المؤمنين صنفين.

\separator

\section{تأويل (ويقيمون الصلاة ومما رزقناهم ينفقون)}
\isnad{قال أبو جعفر:} «وإقامتها إداؤها بحدودها وفروضها والواجب فيها... وأما الصلاة في كلام العرب فإنها الدعاء... وأرى أن الصلاة المفروضة سميت صلاة لأن المصلي متعرض لاستنجاح طلبته من ثواب الله بعمله، مع ما يسأل ربه فيها من حاجاته».

\begin{fawaid}[هل الصلاة في لغة العرب محصورة في الدعاء؟]
تأثر بعض اللغويين والمتكلمين (كالمعتزلة) بفكرة أن الكلمة تُوضع لمعنى ثم تُستعار لغيره، فقالوا إن الصلاة في اللغة هي الدعاء فقط ثم نُقلت شرعاً للعبادة ذات الركوع والسجود. وهذا القول محل نقد؛ فالصلاة معروفة عند العرب كعبادة لها هيئة (كما في قوله: وما كان صلاتهم عند البيت إلا مكاء وتصدية)، وفي أشعارهم. فهي أعم من مجرد الدعاء، وإن كان الدعاء من معانيها الكبرى.
\end{fawaid}

\isnad{وقال في تأويل (ومما رزقناهم ينفقون):} «وأولى التأويلات بالآية... أن يكونوا كانوا مؤدين لجميع اللازم لهم في أموالهم، زكاة كان ذلك أو نفقة من لزمته نفقته من أهل وعيال... لأن الله عمّ وصفهم، ولم يخصص مدحهم بنوع من النفقات».

\separator

\section{تأويل (إن الذين كفروا سواء عليهم)}
\isnad{قال أبو جعفر فيمن نزلت فيهم هذه الآية:} «كان \rawi{ابن عباس} يرى أن هذه الآية نزلت في اليهود الذين كانوا بنواحي المدينة على عهد رسول الله صلى الله عليه وسلم، توبيخاً لهم في جحودهم نبوة محمد صلى الله عليه وسلم وتكذيبهم به مع علمهم به...».

ورجح \rawi{الطبري} هذا القول على من زعم أنها نزلت في قادة الأحزاب الذين قُتلوا ببدر، محتجاً بوحدة السياق والموضوع؛ فبما أن السورة ابتدأت بذكر مؤمني أهل الكتاب، فمن الأليق أن يكون الكلام بعدها عن كفار أهل الكتاب الذين كتموا ما في كتبهم من البينات جحوداً وحسداً، فاستحقوا ألا ينفعهم الإنذار.

\separator

\section{تأويل (ختم الله على قلوبهم) والرد على القدرية}
\isnad{قال أبو جعفر:} «وأصل الختم الطبع... فمعنى الختم عليها وعلى الأسماع التي بها تُدرك المسموعات... نظير معنى الختم على سائر الأوعية والظروف... وقد اختلف أهل التأويل في صفة ذلك».

ثم اختار \rawi{الطبري} لتفسيرها الحديث الحسن: \begin{hadith}إن المؤمن إذا أذنب ذنباً كانت نكتة سوداء في قلبه، فإن تاب ونزع واستغفر صُقل قلبه، وإن زاد زادت حتى تعلو قلبه... فذلك الران\end{hadith}. فأثبت أن الذنوب إذا تتابعت أغلقت القلوب، وأتاها الختم من الله كعقوبة، فلا يكون للإيمان إليها مسلك.

\subsection{رد الطبري على تأويل المعتزلة (القدرية)}
رد \rawi{الطبري} بقوة على المعتزلة الذين زعموا أن معنى (ختم الله) هو إخبار عن "تكبر الكفار وإعراضهم" لينفوا أن يكون الله هو الفاعل لذلك! 

\isnad{فقال ملزماً إياهم:} «فإن زعموا أن ذلك فعل منه (أي من الكافر)، قيل لهم: فإن الله قد أخبر أنه هو الذي ختم على قلوبهم! وكيف يجوز أن يكون إعراض الكافر فعلاً لله؟».

\begin{fawaid}[مسألة تكليف ما لا يطاق]
استدل \rawi{الطبري} بهذه الآية على بطلان قول من أنكر "تكليف ما لا يطاق". وقد جانبه الصواب في هذا الإطلاق؛ فأهل السنة يقررون أن الله لا يكلف نفساً إلا وسعها (أي ما تطيقه من جهة القدرة والأهلية)، أما كون العبد لن يؤمن في علم الله وسابق قدره (كأبي لهب أو هؤلاء المختوم على قلوبهم) فهذا ليس من باب "تكليف ما لا يطاق" الذي يعجز عنه العبد ابتداءً، بل هو من باب العقوبة على الجحود، وتركهم الإيمان كان باختيارهم وكسبهم لا لعدم أهليتهم.
\end{fawaid}

\begin{fawaid}[يحتج الطبري بالسياق في تعيين المخاطبين وتوجيه المعاني]
من المواضع البارزة هنا أن \rawi{الطبري} لا يفسر الجملة منفصلة عن سياق السورة، بل يرجح في أوائل البقرة بأن يربط بين أصناف المؤمنين والكافرين والمنافقين، ويستعمل انتظام السياق من أقوى وجوه الترجيح.
\end{fawaid}

\separator

\section{تأويل (يخادعون الله) والرد على اللغويين}
\isnad{قال أبو جعفر في قوله تعالى:} \begin{ayah}يُخَادِعُونَ اللَّهَ وَالَّذِينَ آمَنُوا وَمَا يَخْدَعُونَ إِلَّا أَنفُسَهُمْ وَمَا يَشْعُرُونَ\end{ayah}: «خداع المنافق ربه والمؤمنين: إظهاره بلسانه من القول والتصديق خلاف الذي في قلبه من الشك والتكذيب ليدرأ عن نفسه... وذلك من فعله وإن كان خداعاً للمؤمنين في عاجل الدنيا فهو لنفسه بذلك خادع... لأنه موردها حياض عطبها».

وقد رد \rawi{الطبري} على \rawi{أبي عبيدة معمر بن المثنى} (من علماء البصرة) الذي ادعى أن (يُخَادِعُونَ) هنا جاءت على صيغة المفاعلة ولكنها بمعنى الفعل من طرف واحد (أي يَخدعون). 

\isnad{فقال الطبري:} «وليس القول في ذلك عندي كالذي قال، بل ذلك من التفاعل الذي لا يكون إلا من اثنين... وذلك أن المنافق يخادع الله بأكاذيبه... والله خادعه بخذلانه عن حسن البصيرة واستدراجه إياه». فأثبت \rawi{الطبري} صفة الخداع لله عز وجل على وجه الكمال والمجازاة للمنافقين.

\separator

بهذا نكون قد أتممنا الجزء الأول من تفسير سورة البقرة، وسنكمل غداً بإذن الله. والسلام عليكم ورحمة الله وبركاته.
